\section{Conclusion}\label{sec:conclusion}

Eight years separate BERT's 80.5 GLUE from Opus~5's perfect IMO. The
progression was not a smooth exponential on one axis but a sequence of
paradigm changes --- fine-tuning, few-shot scale, alignment, reasoning, agency
--- each moving capability onto a new axis and each burning out the previous
era's benchmarks. Our measurements quantify three 2026 facts. Capability in
agentic coding is still compounding at roughly six-fold per year. The cost of
a fixed capability level is falling even faster, to the point where a \$0.20/M
budget tier replicates most of the professionally relevant performance of an
eleven-week-old flagship. And the frontier has fragmented by task: frontend
coding, repository coding, agentic terminal work, and novel reasoning each
have a different leader, which makes routing --- not model selection --- the
decision that matters. The most robust way to improve a deployed system today
is often to change nothing about the model and everything about how its
inference compute is spent.

\paragraph{Reproducibility.} All datasets (\texttt{code/data/}), analysis
scripts (\texttt{exp1}--\texttt{exp4}), the live experiment, unit tests
(\texttt{tests/}), figure generation, and this paper's build script are
included in the accompanying archive; see \texttt{README.md}. Nothing in this
paper requires proprietary access to reproduce except the live benchmark
numbers, whose sources and retrieval dates are documented in
\texttt{code/data\_sources.md}.
